\section{一回の checking cut}

\subsection{制約を解く関係}

$\DDAlignTypes(S,\tau,\rho,S')$ は coercion を伴わない equality alignment，
$\DDAlign(S,\tau,\rho,S')$ は checking cut 全体である．どちらも入力 $S$ を
破壊せず，constraint を解く局所 delta を $S$ の後へ合成して $S'$ を得る．

$\ExactCapMGU$ と $\ExactMGU$ は，それぞれ capability と paired type constraint の
proof-carrying MGU である．MGU の soundness と因子化に加え，constraint 外で恒等，
range が constraint 内，solved form が冪等であることを要求する．この exactness により，
unification と無関係な metavariable を同時に構造化する guess を除く．

$\OneWayDelta(\kappa_p,\tau_p,\kappa_c,\tau_c,D)$ は producer を固定したまま
consumer capability と target を解く exact delta である．consumer $\Any$ は wildcard，
consumer variable は同じ raw occurrence 間で一貫した値を要求する．producer 側の
capability を consumer demand に合わせて強化してはならない．

\subsection{通常の等価整合}

matcher 同士または slot 同士では capability を先に解き，その delta を target へ作用させて
から target constraint を解く．それ以外の型は paired type 全体を一度に解く．
\begin{mathpar}
\inferrule{
  S\tau=\Matcher{\kappa_1}{\tau_1}\quad
  S\rho=\Matcher{\kappa_2}{\tau_2} \\
  \ExactCapMGU(\kappa_1,\kappa_2,C) \\
  \ExactMGU(C\tau_1,C\tau_2,T)
}{
  \DDAlignTypes(S,\tau,\rho,
    \mathsf{seq}(T,\mathsf{seq}((C,\mathsf{id}),S)))
}\rn{ALIGN-MATCHER}

\inferrule{
  S\tau=\MatcherSlot{\kappa_1}{\tau_1}\quad
  S\rho=\MatcherSlot{\kappa_2}{\tau_2} \\
  \ExactCapMGU(\kappa_1,\kappa_2,C) \\
  \ExactMGU(C\tau_1,C\tau_2,T)
}{
  \DDAlignTypes(S,\tau,\rho,
    \mathsf{seq}(T,\mathsf{seq}((C,\mathsf{id}),S)))
}\rn{ALIGN-SLOT}

\inferrule{
  \mathsf{alignPairClass}(S\tau,S\rho)=\mathsf{ordinary} \\
  \ExactMGU(S\tau,S\rho,D)
}{
  \DDAlignTypes(S,\tau,\rho,\mathsf{seq}(D,S))
}\rn{ALIGN-ORDINARY-TYPES}
\end{mathpar}

\subsection{Demand classifier と coercion の選択}

$\dclass(S\tau,S\rho)$ は resolved type だけから次の五つの class を決める．
product of matchers，product of slots の順に調べるため，空 product は
$\mathsf{productMatcherLift}$ が優先される．
\[
\mathsf{productMatcherLift},\quad
\mathsf{slotTupleLift},\quad
\mathsf{matcherToSlot},\quad
\mathsf{slotToSlot},\quad
\mathsf{ordinary}.
\]

\begin{mathpar}
\inferrule{
  \mathsf{productMatcherDuals?}(S\tau)=
    \mathsf{some}\;\overline\eta \\
  S\rho=\MatcherSlot{\kappa}{\upsilon} \\
  \OneWayDelta((\overline{\eta.\kappa}),
    (\overline{\eta.\tau}),\kappa,\upsilon,D)
}{
  \DDAlign(S,\tau,\rho,\mathsf{seq}(D,S))
}\rn{ALIGN-PRODUCT-MATCHER}

\inferrule{
  \dclass(S\tau,S\rho)=\mathsf{slotTupleLift} \\
  \mathsf{productSlotDuals?}(S\tau)=
    \mathsf{some}\;\overline\eta \\
  S\rho=\MatcherSlot{\kappa}{\upsilon} \\
  \ExactCapMGU((\overline{\eta.\kappa}),\kappa,C) \\
  \ExactMGU(C(\overline{\eta.\tau}),C\upsilon,T)
}{
  \DDAlign(S,\tau,\rho,
    \mathsf{seq}(T,\mathsf{seq}((C,\mathsf{id}),S)))
}\rn{ALIGN-SLOT-TUPLE}
\end{mathpar}

\begin{mathpar}
\inferrule{
  S\tau=\Matcher{\kappa_p}{\tau_p}\quad
  S\rho=\MatcherSlot{\kappa_c}{\tau_c} \\
  \OneWayDelta(\kappa_p,\tau_p,\kappa_c,\tau_c,D)
}{
  \DDAlign(S,\tau,\rho,\mathsf{seq}(D,S))
}\rn{ALIGN-MATCHER-SLOT}

\inferrule{
  S\tau=\MatcherSlot{\kappa_1}{\tau_1}\quad
  S\rho=\MatcherSlot{\kappa_2}{\tau_2} \\
  \ExactCapMGU(\kappa_1,\kappa_2,C)\quad
  \ExactMGU(C\tau_1,C\tau_2,T)
}{
  \DDAlign(S,\tau,\rho,
    \mathsf{seq}(T,\mathsf{seq}((C,\mathsf{id}),S)))
}\rn{ALIGN-SLOT-SLOT}

\inferrule{
  \dclass(S\tau,S\rho)=\mathsf{ordinary} \\
  \DDAlignTypes(S,\tau,\rho,S')
}{
  \DDAlign(S,\tau,\rho,S')
}\rn{ALIGN-ORDINARY}
\end{mathpar}

product-of-matchers branch は，whole-product matcher の生成と matcher-to-slot 消費という
二 primitive step を一つの checking cut で行う．別の構文位置でもう一度 coercion する
規則ではない．non-ordinary class はすべて resolved expected head が slot の場合だけである．
matcher-headed expected type は必ず ordinary alignment へ退化する．

\subsection{唯一の checking 規則}

式の checking 規則は次の一つだけである．expected type は synthesis の入力にならず，
synthesis が返したその cut の $S_1$ でだけ alignment を選ぶ．
\begin{mathpar}
\inferrule{
  q;S;\Gamma\vdash e\Rightarrow\tau\dashv q_1;S_1 \\
  \DDAlign(S_1,\tau,\rho,S')
}{
  q;S;\Gamma\vdash e\Leftarrow\rho\dashv q_1;S'
}\rn{DD-CHECK}
\end{mathpar}
